\clearpage
\begin{proofbox}
\textbf{\textcolor{CProof}{先读：三种证明方式各有什么用？}}

下面先围绕核心规律“半周期又出现对称轴/对称中心”，用三种不同方法理解：
\begin{itemize}[leftmargin=*,itemsep=0.25em]
\item \textbf{方法一（门槛最低）：} 数轴上对折、平移，先明白为什么是“一半”。
\item \textbf{方法二（考试最快）：} 把题目给的两个等式直接接起来，几行就结束。
\item \textbf{方法三（最容易举一反三）：} 把图像以对称位置为中间重新写一次，统一处理轴和中心。
\end{itemize}
证明后还补上“两轴推出周期”“全部对称轴”“轴和中心的四分之一周期关系”。

\medskip
\textbf{\textcolor{CProof}{方法一：初中生也能跟上的“折纸 + 数轴”解释}}
\begin{center}\begin{tikzpicture}[font=\small,>=stealth,line join=round]
\draw[gray!40,rounded corners=3pt] (0,0) rectangle (14.2,8.3);\node at (7.1,7.700000000000001) {证明一：右点先移回一个周期，再沿旧轴配对};
\begin{scope}[shift={(2.3857142857142857,3.8222222222222224)},x=0.8571428571428571cm,y=1.277777777777778cm]\draw[->,gray] (-1.5,0)--(5.5,0) node[right]{$x$};\draw[->,gray] (0,-1.7)--(0,1.9) node[above]{$y$};\draw[blue,thick,solid,domain=-1:5,samples=180] plot(\x,{cos(deg(pi*(\x-1)/2))});\fill[orange!85!black] (-0.5,-0.7071067811865476) circle(2.2pt);\fill[orange!85!black] (2.5,-0.7071067811865476) circle(2.2pt);\fill[orange!85!black] (3.5,-0.7071067811865476) circle(2.2pt);\foreach \a in {1,3}{\draw[violet,dashed] (\a,-1.25)--(\a,1.25);\node[above] at (\a,1.3){$x=\a$};}\foreach \X/\L in {-.5/{R-T},2.5/L,3.5/R}{\draw[gray,dashed] (\X,-.70710678)--(\X,-1.2);\node[below=4pt] at (\X,-1.4){$\L$};}\end{scope}\node[anchor=west,align=left,text width=5.0cm] at (8.4,4.0) {$T=4,\ a=1,\ t=\frac12$\\[9pt]新轴 $a+T/2=3$\\[9pt]$R=3.5,\ R-T=-0.5$\\[9pt]$L=2.5$\\[9pt]三处都为 $-\frac{\sqrt2}2$};
\node[text width=12.6cm,align=center] at (7.1,.65) {图取一组具体距离；证明中t任意，平移一次与沿旧轴配对就得到新轴。};
\end{tikzpicture}\end{center}



设 $x=a$ 是一条对称轴，图像每走 $T$ 个单位就重复。

\textbf{第一步。} 取新候选折痕 $a+T/2$。在它左右各取距离 $t$ 的点：
\[
L=a+\frac T2-t,\qquad R=a+\frac T2+t.
\]

\textbf{第二步。} 把右点向左搬一个完整周期：
\[
R-T=a-\frac T2+t.
\]

\textbf{第三步。} 看两个点 $R-T$ 与 $L$：
\[
\frac{(R-T)+L}{2}=a.
\]
所以它们关于旧折痕 $x=a$ 对称；根据题目条件，两处函数值相同：
\[
f(R-T)=f(L).
\]
又因为 $R$ 和 $R-T$ 相隔 $T$，函数值也相同，故
\[
f(R)=f(R-T)=f(L).
\]

\textbf{第四步。} $t$ 可以随便取，所以新折痕的任意左右两点函数值都相同。\textbf{这就证明 $x=a+T/2$ 也是一条对称轴。}

\textbf{对称中心怎么理解？} 在 $(b,c)$ 左右同样远的两点，函数值\textbf{相加是 $2c$}；把右边的点搬回一个完整周期，正好落到旧中心另一侧。因此在新中心 $(b+T/2,c)$ 的左右两点，函数值之和也仍是 $2c$。与证明对称轴的步骤完全平行，只是“相等”换成了“相加等于 $2c$”。

\medskip

\end{proofbox}
\clearpage
\begin{proofbox}
\textbf{\textcolor{CProof}{方法二：考场上最快的直接代换}}
\begin{center}\begin{tikzpicture}[font=\small,>=stealth,line join=round]
\draw[gray!40,rounded corners=3pt] (0,0) rectangle (14.2,8.3);\node at (7.1,7.700000000000001) {例3：中心横向移3，高度始终保留5};
\begin{scope}[shift={(1.4000000000000001,1.8921052631578945)},x=0.6cm,y=0.4842105263157894cm]\draw[->,gray] (-0.5,0)--(9.5,0) node[right]{$x$};\draw[->,gray] (0,-0.5)--(0,9) node[above]{$y$};\draw[teal,thick,solid,domain=0:9,samples=180] plot(\x,{5+2*sin(deg(pi*(\x-2)/3))});\fill[orange!85!black] (2,5) circle(2.2pt);\fill[orange!85!black] (5,5) circle(2.2pt);\fill[orange!85!black] (8,5) circle(2.2pt);\fill[orange!85!black] (3,6.732050807568877) circle(2.2pt);\fill[orange!85!black] (7,3.267949192431123) circle(2.2pt);\draw[gray,dashed] (0,5)--(9,5);\foreach \v in {2,5,8}{\node[below=4pt] at (\v,0){$\v$};}\end{scope}\node[anchor=west,align=left,text width=5.0cm] at (8.4,4.0) {$f=5+2\sin\frac{\pi(x-2)}3$\\[9pt]$T=6$，中心 $(2,5),(5,5)$\\[9pt]$f(5)=f(8)=5$\\[9pt]$f(3)+f(7)=10$\\[9pt]纵坐标不增加 $T/2$};
\node[text width=12.6cm,align=center] at (7.1,.65) {橙点包括中心与绕中心配对的点；相邻中心横坐标相差3，高度都为5。};
\end{tikzpicture}\end{center}



\textbf{（1）证明对称轴。} 已知
\[
f(x+T)=f(x),\qquad f(a+t)=f(a-t).
\]
只要证明对任意 $t$，
\[
f\left(a+\frac T2+t\right)=f\left(a+\frac T2-t\right).
\]
逐步推算：
\[
\begin{aligned}
f\left(a+\frac T2+t\right)
&=f\left(a-\frac T2+t\right) &&\text{（左移 $T$）}\\
&=f\left(a-\left(\frac T2-t\right)\right)\\
&=f\left(a+\left(\frac T2-t\right)\right) &&\text{（沿 $x=a$ 对称）}\\
&=f\left(a+\frac T2-t\right).
\end{aligned}
\]
证毕。向左也可同样推，反复做便得 $x=a+kT/2$。

\textbf{（2）证明对称中心。} 已知
\[
f(b+t)+f(b-t)=2c.
\]
任取 $t$，有
\[
\begin{aligned}
&f\left(b+\frac T2+t\right)+f\left(b+\frac T2-t\right)\\
&\qquad=f\left(b-\frac T2+t\right)+f\left(b+\frac T2-t\right)
&&\text{（第一项左移 $T$）}\\
&\qquad=f\left(b-\left(\frac T2-t\right)\right)
  +f\left(b+\left(\frac T2-t\right)\right)\\
&\qquad=2c &&\text{（旧中心的对称关系）}.
\end{aligned}
\]
故 $(b+T/2,c)$ 也是对称中心。反复向左右移半周期，就得到全部\emph{保证存在}的这一组中心。

\medskip

\end{proofbox}
\clearpage
\begin{proofbox}
\textbf{\textcolor{CProof}{方法三：先把对称位置移到原点，一次看懂两类对称}}
\begin{center}\begin{tikzpicture}[font=\small,>=stealth,line join=round]
\draw[gray!40,rounded corners=3pt] (0,0) rectangle (14.2,8.8);\node at (7.1,8.200000000000001) {证明三：把中心移到原点，转成奇函数关系};
\node at (3.1500000000000004,7.3) {$f(x)=3+\sin\pi(x-\frac12)$};\begin{scope}[shift={(1.6800000000000002,3.4499999999999997)},x=0.9800000000000001cm,y=0.5937499999999999cm]\draw[->,gray] (-1,0)--(4,0) node[right]{$x$};\draw[->,gray] (0,-1.6)--(0,4.8) node[above]{$y$};\draw[teal,thick,solid,domain=-0.7:3.7,samples=180] plot(\x,{3+sin(deg(pi*(\x-.5)))});\fill[orange!85!black] (0.5,3) circle(2pt);\fill[orange!85!black] (1.5,3) circle(2pt);\draw[gray,dashed] (-.8,3)--(3.8,3);\end{scope}\node[align=center,text width=5.8cm] at (3.1500000000000004,1.7) {中心 $(\frac12,3)$\\原来的周期为2};\node at (9.95,7.3) {$h(t)=f(\frac12+t)-3$};\begin{scope}[shift={(8.48,3.4499999999999997)},x=0.9800000000000001cm,y=0.5937499999999999cm]\draw[->,gray] (-1,0)--(4,0) node[right]{$t$};\draw[->,gray] (0,-1.6)--(0,4.8) node[above]{$y$};\draw[orange!85!black,thick,solid,domain=-0.7:3.7,samples=180] plot(\x,{sin(deg(pi*\x))});\fill[orange!85!black] (0,0) circle(2pt);\fill[orange!85!black] (1,0) circle(2pt);\end{scope}\node[align=center,text width=5.8cm] at (9.95,1.7) {中心变为 $(0,0)$\\$h=\sin\pi t$，周期仍为2};
\node[text width=12.6cm,align=center] at (7.1,.65) {两图同尺度；坐标替换和减去中心高度，不改变横向重复的距离。};
\end{tikzpicture}\end{center}



\textbf{对称轴的情况。} 设 $g(t)=f(a+t)$。旧对称轴是 $x=a$，所以
\[
g(-t)=g(t).
\]
函数有周期 $T$，所以 $g(t+T)=g(t)$。于是
\[
\begin{aligned}
g\left(\frac T2+t\right)
&=g\left(-\frac T2+t\right) &&\text{（相隔 $T$）}\\
&=g\left(\frac T2-t\right) &&\text{（左右相反时值相同）}.
\end{aligned}
\]
这说明 $g$ 在 $T/2$ 左右也对称，换回 $f$ 就得到新对称轴 $x=a+T/2$。

\textbf{对称中心的情况。} 设 $h(t)=f(b+t)-c$，也就是先让中心横坐标变成 $0$，再把中心高度减到 $0$。于是
\[
h(-t)=-h(t),\qquad h(t+T)=h(t).
\]
同样算：
\[
\begin{aligned}
h\left(\frac T2+t\right)
&=h\left(-\frac T2+t\right)\\
&=-h\left(\frac T2-t\right).
\end{aligned}
\]
因此
\[
h\left(\frac T2+t\right)+h\left(\frac T2-t\right)=0.
\]
把减掉的 $c$ 加回去，立刻得到新对称中心 $(b+T/2,c)$。

\medskip
\end{proofbox}\clearpage\begin{proofbox}\textbf{\textcolor{CProof}{再证结论三：为什么两条对称轴之间有“两倍间距”的周期？}}
\begin{center}\begin{tikzpicture}[font=\small,>=stealth,line join=round]
\draw[gray!40,rounded corners=3pt] (0,0) rectangle (14.2,8.3);\node at (7.1,7.700000000000001) {例4：关于负1与3反射两次，横向总共移动8};
\begin{scope}[shift={(2.7153846153846155,3.8222222222222224)},x=0.46153846153846156cm,y=1.277777777777778cm]\draw[->,gray] (-3.5,0)--(9.5,0) node[right]{$x$};\draw[->,gray] (0,-1.7)--(0,1.9) node[above]{$y$};\draw[blue,thick,solid,domain=-3:9,samples=180] plot(\x,{cos(deg(pi*(\x+1)/4))});\fill[orange!85!black] (0,0.7071067811865476) circle(2.2pt);\fill[orange!85!black] (-2,0.7071067811865476) circle(2.2pt);\fill[orange!85!black] (8,0.7071067811865476) circle(2.2pt);\foreach \a in {-1,3}{\draw[violet,dashed] (\a,-1.3)--(\a,1.3);\node[above] at (\a,1.35){$x=\a$};}\end{scope}\node[anchor=west,align=left,text width=5.0cm] at (8.4,4.0) {$f=\cos\frac{\pi(x+1)}4$\\[9pt]轴 $x=-1,3$\\[9pt]$0\to-2\to8$\\[9pt]三点高度 $\frac{\sqrt2}2$\\[9pt]周期可取 $2\times4=8$};
\node[text width=12.6cm,align=center] at (7.1,.65) {两次反射的合成位移，是两轴间距的两倍；这里只证明一个周期。};
\end{tikzpicture}\end{center}



旧对称轴 $x=a_1$ 保证：对任意 $u$，
\[
f(u)=f(2a_1-u).
\]
另一条对称轴 $x=a_2$ 再保证：
\[
f(2a_1-u)=f\bigl(2a_2-(2a_1-u)\bigr).
\]
合在一起便是
\[
\boxed{f(u)=f\bigl(u+2(a_2-a_1)\bigr)}.
\]
如果 $a_2>a_1$，这表示向右移动 $2(a_2-a_1)$ 后函数值不变。如果 $a_2<a_1$，则向左移动 $2(a_1-a_2)$ 不变，换个方向同样可以得到一个正周期。所以 $2|a_2-a_1|$ 是一个正周期。

\medskip
\end{proofbox}\clearpage\begin{proofbox}\textbf{\textcolor{CProof}{为什么知道最小正周期以后，可以列出全部对称轴？}}
\begin{center}\begin{tikzpicture}[font=\small,>=stealth,line join=round]
\draw[gray!40,rounded corners=3pt] (0,0) rectangle (14.2,8.3);\node at (7.1,7.700000000000001) {最小正周期为2时，全部轴恰好位于整数处};
\begin{scope}[shift={(2.3,3.8222222222222224)},x=1.2cm,y=1.277777777777778cm]\draw[->,gray] (-1,0)--(4,0) node[right]{$x$};\draw[->,gray] (0,-1.7)--(0,1.9) node[above]{$y$};\draw[blue,thick,solid,domain=-0.7:3.7,samples=180] plot(\x,{cos(deg(pi*\x))});\fill[orange!85!black] (0.2,0.8090169943749475) circle(2.2pt);\fill[orange!85!black] (0.6,-0.30901699437494734) circle(2.2pt);\foreach \a in {0,1,2,3}{\draw[violet,dashed] (\a,-1.3)--(\a,1.3);}\draw[red,dotted,thick] (.4,-1.3)--(.4,1.3);\end{scope}\node[anchor=west,align=left,text width=5.0cm] at (8.4,4.0) {$f=\cos\pi x,\ T_0=2$\\[9pt]全部轴：$x=k$\\[9pt]候选 $x=0.4$ 不成立\\[9pt]$f(0.2)\ne f(0.6)$\\[9pt]非整数处不能添一条轴};
\node[text width=12.6cm,align=center] at (7.1,.65) {图只示范排除一个错误候选；正文用最小正周期的整数倍性质排除所有其他位置。};
\end{tikzpicture}\end{center}



设最小正周期是 $T_0$，已有一条对称轴 $x=a$。

\textbf{先证明“这些都是”。} 由上面的方法，$x=a+kT_0/2$ 都是对称轴。

\textbf{再证明“不会有别的”。} 假设 $x=s$ 是另一条对称轴，那么两轴推出
\[
2|s-a|\quad\text{是一个正周期}.
\]
\textbf{任何正周期 $P$ 都必须是 $T_0$ 的整数倍。} 道理很简单：假如 $P=nT_0+r$，其中 $n$ 是非负整数，$0\le r<T_0$，那么“先走 $n$ 个最小周期，再把剩下的 $r$ 走完”也应该重复；因此剩下的 $r$ 本身也是周期。若 $r>0$ 就会出现一个比最小正周期更小的正周期，矛盾，所以只能有 $r=0$。

因此 $2|s-a|=nT_0$。换上正负方向，可以写成
\[
\boxed{s=a+\frac{kT_0}{2},\quad k\in\mathbb Z}.
\]
这就把全部对称轴找齐了。注意：\textbf{只有已知 $T_0$ 真的是最小正周期，这一步才成立。}

\medskip
\end{proofbox}\clearpage\begin{proofbox}\textbf{\textcolor{CProof}{进阶但仍属高中范围：为什么是“四分之一周期”？}}
\begin{center}\begin{tikzpicture}[font=\small,>=stealth,line join=round]
\draw[gray!40,rounded corners=3pt] (0,0) rectangle (14.2,8.3);\node at (7.1,7.700000000000001) {轴到中心不一定最近：也可能跨3个四分之一周期};
\begin{scope}[shift={(2.158823529411765,3.8222222222222224)},x=1.7647058823529411cm,y=1.277777777777778cm]\draw[->,gray] (-0.6,0)--(2.8,0) node[right]{$x$};\draw[->,gray] (0,-1.7)--(0,1.9) node[above]{$y$};\draw[blue,thick,solid,domain=-0.3:2.5,samples=180] plot(\x,{cos(deg(pi*\x))});\fill[orange!85!black] (0,1) circle(2.2pt);\fill[orange!85!black] (0.5,0) circle(2.2pt);\fill[orange!85!black] (1,-1) circle(2.2pt);\fill[orange!85!black] (1.5,0) circle(2.2pt);\draw[violet,dashed] (0,-1.3)--(0,1.3);\foreach \v/\L in {.5/\frac12,1/1,1.5/\frac32,2/2}{\node[below=4pt] at (\v,-1.35){$\L$};}\end{scope}\node[anchor=west,align=left,text width=5.0cm] at (8.4,4.0) {$T_0=2$，轴 $x=0$\\[9pt]中心 $(\frac12,0),(\frac32,0)$\\[9pt]$\frac12=T_0/4$\\[9pt]$\frac32=3T_0/4$\\[9pt]一般为奇数个 $T_0/4$};
\node[text width=12.6cm,align=center] at (7.1,.65) {四分之一周期关系使用最小正周期；不能把较远中心也误当最近的中心。};
\end{tikzpicture}\end{center}



设 $x=a$ 是对称轴，$(b,c)$ 是对称中心。先沿 $x=a$ 折一次：
\[
f(x)=f(2a-x).
\]
再利用绕 $(b,c)$ 转半圈：
\[
f(2a-x)+f\bigl(2b-(2a-x)\bigr)=2c.
\]
整理得到
\[
\boxed{f\bigl(x+2(b-a)\bigr)=2c-f(x).}\tag{$*$}
\]
因此移动距离 $2(b-a)$，不是一定回到原值，而是把函数值绕水平线 $y=c$ 上下交换。再移动一次：
\[
\begin{aligned}
f\bigl(x+4(b-a)\bigr)
&=2c-f\bigl(x+2(b-a)\bigr)\\
&=2c-[2c-f(x)]\\
&=f(x).
\end{aligned}
\]
如果 $a\ne b$，那么 $4|b-a|$ 就是一个正周期。

若已知 $T_0$ 是最小正周期，前面已证“任何正周期都是它的整数倍”，故
\[
4|b-a|=nT_0\qquad(n=1,2,3,\ldots).
\]
\textbf{这里的 $n$ 还必须是奇数。} 如果 $n$ 是偶数，则 $2(b-a)$ 正好是 $T_0$ 的整数倍，于是左边 $(*)$ 就等于 $f(x)$，得到 $f(x)=2c-f(x)$，也就是 $f(x)=c$ 对所有 $x$ 成立。这样函数是常数函数，与“存在最小正周期”矛盾。于是
\[
\boxed{b-a=\frac{(2k+1)T_0}{4}\quad(k\in\mathbb Z).}
\]
这里的正负号由“中心在轴的左边还是右边”决定。

\textbf{特殊情况：}如果 $a=b$，直接把 $(*)$ 写成 $f(x)=2c-f(x)$，得到 $f(x)=c$。所以\textbf{非恒定函数的同一条竖直线上，不可能既有对称轴又有对称中心。}

\medskip
\end{proofbox}\clearpage\begin{proofbox}\textbf{\textcolor{CProof}{补充辨析：什么时候也能找齐全部对称中心？}}
\begin{center}\begin{tikzpicture}[font=\small,>=stealth,line join=round]
\draw[gray!40,rounded corners=3pt] (0,0) rectangle (14.2,8.8);\node at (7.1,8.200000000000001) {全部中心的补充证明：有固定上下界，才排得掉高度漂移};
\node at (3.1500000000000004,7.3) {有界：$f=\cos\pi x$};\begin{scope}[shift={(1.5166666666666666,3.45)},x=0.8166666666666668cm,y=0.475cm]\draw[->,gray] (-1,0)--(5,0) node[right]{$x$};\draw[->,gray] (0,-2)--(0,6) node[above]{$y$};\draw[blue,thick,solid,domain=-0.7:5,samples=180] plot(\x,{cos(deg(pi*\x))});\fill[orange!85!black] (0.5,0) circle(2pt);\fill[orange!85!black] (1.5,0) circle(2pt);\fill[orange!85!black] (2.5,0) circle(2pt);\draw[gray,dashed] (-.8,1)--(4.8,1);\draw[gray,dashed] (-.8,-1)--(4.8,-1);\end{scope}\node[align=center,text width=5.8cm] at (3.1500000000000004,1.7) {所有值始终在 $[-1,1]$\\中心的高度不能不断增加};\node at (9.95,7.3) {漂移示例：$g=x$};\begin{scope}[shift={(8.316666666666666,3.45)},x=0.8166666666666668cm,y=0.475cm]\draw[->,gray] (-1,0)--(5,0) node[right]{$x$};\draw[->,gray] (0,-2)--(0,6) node[above]{$y$};\draw[orange!85!black,thick,solid,domain=-0.7:5,samples=180] plot(\x,{\x});\fill[orange!85!black] (0,0) circle(2pt);\fill[orange!85!black] (2,2) circle(2pt);\fill[orange!85!black] (1,1) circle(2pt);\fill[orange!85!black] (5,5) circle(2pt);\end{scope}\node[align=center,text width=5.8cm] at (9.95,1.7) {中心 $(0,0),(2,2)$\\迭代后高度从1变5再变9};
\node[text width=12.6cm,align=center] at (7.1,.65) {右图不是周期函数；有界性排除高度漂移。};
\end{tikzpicture}\end{center}



对最一般的抽象函数，不应仅凭最小正周期 $T_0$ 就随意说“全部对称中心”只有 $(b+kT_0/2,c)$。但若额外知道\textbf{函数值始终有一个固定的上下界}，就能证明它们确实是全部。

证明也不难。假设 $(b,c)$ 与 $(d,e)$ 都是对称中心，那么连续使用两次“绕中心转半圈”的等式：
\[
\begin{aligned}
f\bigl(x+2(d-b)\bigr)
&=2e-f(2b-x)\\
&=2e-[2c-f(x)]\\
&=f(x)+2(e-c).
\end{aligned}
\]
若 $e\ne c$，每多移动 $2(d-b)$，函数值就会不断增减同一个非零数；反复移动，函数值将越来越大或越来越小，与“始终有上下界”矛盾。所以 $e=c$，此时 $2(d-b)$ 才是一个真正的周期。

若有最小正周期 $T_0$，便得 $d-b=kT_0/2$。因此\textbf{当函数值有上下界时}，若给定一个对称中心 $(b,c)$，则全部对称中心恰为
\[
\boxed{\left(b+\frac{kT_0}{2},\,c\right),\quad k\in\mathbb Z.}
\]
这条补充适用于常见的正弦、余弦函数；它\textbf{不是}原先无需附加条件就能套用的结论。
\end{proofbox}
